% !TEX root = main.tex

%%%%%%%%%%%%%%%%%%%%%%%%%%%%%%%%%%%
%%% Discussion / Ablation study %%%
%%%%%%%%%%%%%%%%%%%%%%%%%%%%%%%%%%%

\section{Discussion}

The current standard for the evaluation of an unsupervised method involves the use of an AlexNet architecture trained on ImageNet and tested on class-level tasks.
To understand and measure the various biases introduced by this pipeline on \OURS, we consider a different training set, a different architecture and an instance-level recognition task.

\begin{table}[t!]
  \centering
  \begin{tabular}{@{}lc cc cc c cc c cc@{}}
    \toprule
                      &~~~~&            &~~~~& \multicolumn{2}{c}{Classification} &~~& \multicolumn{2}{c}{Detection} &~~& \multicolumn{2}{c}{Segmentation} \\
                      \cmidrule{5-6} \cmidrule{8-9} \cmidrule{11-12}
    Method    && Training set && \textsc{fc6-8}& \textsc{all}   && \textsc{fc6-8}& \textsc{all}  && \textsc{fc6-8}& \textsc{all} \\
                      \midrule
    Best competitor &&ImageNet&&  $~63.0~$ & $~67.7~$ && $~43.4^{\dagger}$ & $~53.2~$ && $~35.8^{\dagger}$ & $~37.7~$ \\
                      \midrule
    \OURS&& ImageNet   && $~72.0~$  & $~73.7~$  && $~51.4~$ & $~55.4~$ && $~43.2~$ & $~45.1~$ \\
    \OURS&& YFCC100M   && $~67.3~$  & $~69.3~$  && $~45.6~$ & $~53.0~$ && $~39.2~$ & $~42.2~$ \\
    \bottomrule
  \end{tabular}
  \vspace{\soustable}
  \caption{
    Impact of the training set on the performance of \OURS measured on the \textsc{Pascal} VOC transfer tasks as described in Sec.~\ref{sec:pascal}.
    We compare ImageNet with a subset of $1$M images from YFCC100M~\cite{thomee2015new}.
    Regardless of the training set, \OURS outperforms the best published numbers on most tasks.
  Numbers for other  methods produced by us are marked with a ${\dagger}$
  }
  \label{tab:flickr}
\end{table}

\subsection{ImageNet versus YFCC100M}

ImageNet is a dataset designed for a fine-grained object classification challenge~\cite{russakovsky2015imagenet}.
It is object oriented, manually annotated and organised into well balanced object categories.
By design, \OURS favors balanced clusters and, as discussed above, our number of cluster $k$ is somewhat comparable with the number of labels in ImageNet.
This may have given an unfair advantage to \OURS over other unsupervised approaches when trained on ImageNet.
To measure the impact of this effect, we consider a subset of randomly-selected $1$M images from the YFCC100M dataset~\cite{thomee2015new} for the pre-training.
Statistics on the hashtags used in YFCC100M suggests that the underlying ``object classes'' are severly unbalanced~\cite{joulin2016learning},
leading to a data distribution less favorable to \OURS.

Table~\ref{tab:flickr} shows the difference in performance on \textsc{Pascal} VOC of \OURS pre-trained on YFCC100M compared to ImageNet.
As noted by Doersch~\etal~\cite{doersch2015unsupervised}, this dataset is not object oriented, hence the performance are expected to drop by a few percents.
However, even when trained on uncured Flickr images, \OURS outperforms the current state of the art by a significant margin on most tasks (up to $+4.3\%$ on classification and $+4.5\%$ on semantic segmentation).
We report the rest of the results in the supplementary material with similar conclusions.
This experiment validates that \OURS is robust to a change of image distribution, leading
to state-of-the-art general-purpose visual features even if this distribution is not favorable to its design.

\subsection{AlexNet versus VGG}

In the supervised setting, deeper architectures like VGG or ResNet~\cite{he2015delving} have a much higher accuracy on ImageNet than AlexNet.
We should expect the same improvement if these architectures are used with an unsupervised approach.
Table~\ref{tab:vocVGG} compares a VGG-16 and an AlexNet trained with \OURS on ImageNet and tested on the \textsc{Pascal} VOC 2007 object detection task with fine-tuning.
We also report the numbers obtained with other unsupervised approaches~\cite{doersch2015unsupervised,wang2017transitive}.
Regardless of the approach, a deeper architecture leads to a significant improvement in performance on the target task.
Training the VGG-16 with \OURS gives a performance above the state of the art, bringing us to only $1.4$ percents below the supervised topline.
Note that the difference between unsupervised and supervised approaches remains in the same ballpark for both architectures (\ie $1.4\%$).
Finally, the gap with a random baseline grows for larger architectures,
justifying the relevance of unsupervised pre-training for complex architectures when little supervised data is available.

\begin{table}[t!]
  \parbox[c][16em][t]{.45\linewidth}{
    \centering
    \begin{tabular}{@{}l c c c c@{}}
      \toprule
      Method &~~~& AlexNet && VGG-16 \\
      \midrule
      ImageNet labels                 && $56.8$ && $67.3$ \\
      Random                          && $47.8$ && $39.7$ \\
      \midrule
      Doersch~\etal~\cite{doersch2015unsupervised}  && $51.1$ && $61.5$ \\
      Wang and Gupta~\cite{wang2015unsupervised}    && $47.2$ && $60.2$ \\
      Wang~\etal~\cite{wang2017transitive}          && --   && $63.2$ \\
      \midrule
      \OURS                           && $\textbf{55.4}$ && $\textbf{65.9}$ \\
      \bottomrule
    \end{tabular}
    \vspace{\soustable}
    \caption{
      \textsc{Pascal} VOC 2007 object detection with AlexNet and VGG-16.
      Numbers are taken from Wang~\etal~\cite{wang2017transitive}.
    }
    \label{tab:vocVGG}
  }
  \hfill
  \parbox[c][16em][t]{.45\linewidth}{
    \centering
    \begin{tabular}{@{}lc c c c@{}}
      \toprule
      Method &~~~& Oxford$5$K && Paris$6$K \\
      \midrule
      ImageNet labels              && $72.4$ && $81.5$  \\
      Random                       && $~6.9$  && $22.0$ \\
      \midrule
      Doersch~\etal~\cite{doersch2015unsupervised}  && $35.4$ && $53.1$ \\
      Wang~\etal~\cite{wang2017transitive}          && $42.3$ && $58.0$ \\
      \midrule
      \OURS                           && $\textbf{61.0}$  && $\textbf{72.0}$ \\
      \bottomrule
    \end{tabular}
    \vspace{\soustable}
    \caption{
      mAP on instance-level image retrieval on Oxford and Paris dataset with a VGG-16.
      We apply R-MAC with a resolution of $1024$ pixels and $3$ grid levels~\cite{tolias2015particular}.
    }
    \label{tab:retrieval}
  }
\end{table}

\subsection{Evaluation on instance retrieval}

The previous benchmarks measure the capability of an unsupervised network to capture class level information.
They do not evaluate if it can differentiate images at the instance level.
To that end, we propose image retrieval as a down-stream task.
We follow the experimental protocol of Tolias~\etal~\cite{tolias2015particular} on two datasets, \ie, Oxford Buildings~\cite{Philbin07} and Paris~\cite{Philbin08}.
Table~\ref{tab:retrieval} reports the performance of a VGG-16 trained with different approaches obtained with Sobel filtering, except for Doersch~\etal~\cite{doersch2015unsupervised} and Wang~\etal~\cite{wang2017transitive}.
This preprocessing improves by $5.5$ points the mAP of a supervised VGG-16 on the Oxford dataset, but not on Paris.
This may translate in a similar advantage for \OURS, but it does not account for the average differences of $19$ points.
Interestingly, random convnets perform particularly poorly on this task compared to pre-trained models.
This suggests that image retrieval is a task where the pre-training is essential and studying it as a down-stream task could give further insights about the quality of the features produced by unsupervised approaches.
